\documentclass[10pt,a4paper]{amsart}
\usepackage{amsmath,amssymb,amsthm,graphicx}
% Retained preamble for the Stacks Project chapter snapshot in this directory.
% The chapter's own preamble.tex is a copy of the shared Stacks Project preamble; the block
% below is sliced MECHANICALLY, line for line, from it. Nothing was added, removed or
% redefined; the wrapper class, the xr-hyper/hyperref wiring and the externaldocument list
% are omitted because the assembler supplies the document class.
%
% lines taken from preamble.tex: 3, 5, 9, 11, 12, 156, 157, 158, ... (51 lines total)
%
% package coverage restored from the chapter's own preamble (2026-10-13): xy, xy-2cell

\usepackage[all]{xy}
\xyoption{2cell}
\UseAllTwocells

\newenvironment{reference}{\comment}{\endcomment}
\newenvironment{slogan}{\comment}{\endcomment}
\newenvironment{history}{\comment}{\endcomment}
\theoremstyle{plain}
\newtheorem{theorem}[subsection]{Theorem}
\newtheorem{proposition}[subsection]{Proposition}
\newtheorem{lemma}[subsection]{Lemma}
\theoremstyle{definition}
\newtheorem{definition}[subsection]{Definition}
\newtheorem{example}[subsection]{Example}
\newtheorem{exercise}[subsection]{Exercise}
\newtheorem{situation}[subsection]{Situation}
\theoremstyle{remark}
\newtheorem{remark}[subsection]{Remark}
\newtheorem{remarks}[subsection]{Remarks}
\numberwithin{equation}{subsection}
\def\lim{\mathop{\mathrm{lim}}\nolimits}
\def\colim{\mathop{\mathrm{colim}}\nolimits}
\def\Spec{\mathop{\mathrm{Spec}}}
\def\Hom{\mathop{\mathrm{Hom}}\nolimits}
\def\Ext{\mathop{\mathrm{Ext}}\nolimits}
\def\SheafHom{\mathop{\mathcal{H}\!\mathit{om}}\nolimits}
\def\SheafExt{\mathop{\mathcal{E}\!\mathit{xt}}\nolimits}
\def\Sch{\mathit{Sch}}
\def\Mor{\mathop{\mathrm{Mor}}\nolimits}
\def\Ob{\mathop{\mathrm{Ob}}\nolimits}
\def\Sh{\mathop{\mathit{Sh}}\nolimits}
\def\NL{\mathop{N\!L}\nolimits}
\def\CH{\mathop{\mathrm{CH}}\nolimits}
\def\proetale{{pro\text{-}\acute{e}tale}}
\def\etale{{\acute{e}tale}}
\def\QCoh{\mathit{QCoh}}
\def\Ker{\mathop{\mathrm{Ker}}}
\def\Im{\mathop{\mathrm{Im}}}
\def\Coker{\mathop{\mathrm{Coker}}}
\def\Coim{\mathop{\mathrm{Coim}}}
\DeclareMathSymbol{\boxtimes}{\mathbin}{AMSa}{"02}
\def\QCohstack{\mathcal{QC}\!\mathit{oh}}
\def\Cohstack{\mathcal{C}\!\mathit{oh}}
\def\Spacesstack{\mathcal{S}\!\mathit{paces}}
\def\Quotfunctor{\mathrm{Quot}}
\def\Hilbfunctor{\mathrm{Hilb}}
\def\Curvesstack{\mathcal{C}\!\mathit{urves}}
\def\Polarizedstack{\mathcal{P}\!\mathit{olarized}}
\def\Complexesstack{\mathcal{C}\!\mathit{omplexes}}
\def\Pic{\mathop{\mathrm{Pic}}\nolimits}
\def\Picardstack{\mathcal{P}\!\mathit{ic}}
\def\Picardfunctor{\mathrm{Pic}}
\def\Deformationcategory{\mathcal{D}\!\mathit{ef}}
\usepackage{cleveref}
\crefname{originallemma}{original lemma}{original lemmas}
\Crefname{originallemma}{Original lemma}{Original lemmas}
\makeatletter
\newlabel{algebra-lemma-NL-homotopy}{{00S1}{source}{}{sourcepointer-1}{}}
\newlabel{algebra-lemma-NL-homotopy@cref}{{[originallemma][0][]00S1}{[1][0][]source}{}{}{}}
\makeatother
\crefname{originallemma}{original lemma}{original lemmas}
\Crefname{originallemma}{Original lemma}{Original lemmas}
\makeatletter
\newlabel{algebra-section-netherlander}{{00S0}{source}{}{sourcepointer-2}{}}
\newlabel{algebra-section-netherlander@cref}{{[originallemma][0][]00S0}{[1][0][]source}{}{}{}}
\makeatother
\crefname{originallemma}{original lemma}{original lemmas}
\Crefname{originallemma}{Original lemma}{Original lemmas}
\makeatletter
\newlabel{lemma-first-order-completion}{{0DZ4}{source}{}{sourcepointer-3}{}}
\newlabel{lemma-first-order-completion@cref}{{[originallemma][0][]0DZ4}{[1][0][]source}{}{}{}}
\makeatother
\crefname{originallemma}{original lemma}{original lemmas}
\Crefname{originallemma}{Original lemma}{Original lemmas}
\makeatletter
\newlabel{lemma-henselization}{{0DYT}{source}{}{sourcepointer-4}{}}
\newlabel{lemma-henselization@cref}{{[originallemma][0][]0DYT}{[1][0][]source}{}{}{}}
\makeatother
\crefname{originallemma}{original lemma}{original lemmas}
\Crefname{originallemma}{Original lemma}{Original lemmas}
\makeatletter
\newlabel{more-algebra-lemma-base-change-RHom}{{0A6A}{source}{}{sourcepointer-5}{}}
\newlabel{more-algebra-lemma-base-change-RHom@cref}{{[originallemma][0][]0A6A}{[1][0][]source}{}{}{}}
\makeatother
\crefname{originallemma}{original lemma}{original lemmas}
\Crefname{originallemma}{Original lemma}{Original lemmas}
\makeatletter
\newlabel{more-algebra-lemma-henselization-NL}{{0D08}{source}{}{sourcepointer-6}{}}
\newlabel{more-algebra-lemma-henselization-NL@cref}{{[originallemma][0][]0D08}{[1][0][]source}{}{}{}}
\makeatother
\crefname{originallemma}{original lemma}{original lemmas}
\Crefname{originallemma}{Original lemma}{Original lemmas}
\makeatletter
\newlabel{more-algebra-lemma-henselization-flat}{{0AGU}{source}{}{sourcepointer-7}{}}
\newlabel{more-algebra-lemma-henselization-flat@cref}{{[originallemma][0][]0AGU}{[1][0][]source}{}{}{}}
\makeatother
\crefname{originallemma}{original lemma}{original lemmas}
\Crefname{originallemma}{Original lemma}{Original lemmas}
\makeatletter
\newlabel{more-algebra-lemma-neighbourhood-equivalence}{{05EC}{source}{}{sourcepointer-8}{}}
\newlabel{more-algebra-lemma-neighbourhood-equivalence@cref}{{[originallemma][0][]05EC}{[1][0][]source}{}{}{}}
\makeatother
\title{Stacks, Examples of Deformations: lemma-first-order-henselization}
\author{The Stacks Project Authors}
\date{stacks-project:examples-defos.tex:local-snapshot-20261012; GNU FDL 1.2 or later}
\begin{document}
\maketitle
\noindent\textit{Editorial note.} Statement span is lines 2773-2785 of sources/stacks-examples-defos/pinned-original-source.tex; proof span is lines 2787-2832, the complete author proof (44 body lines).\par\medskip
\noindent\textit{Editorial note.} Navigation stubs: 8 cited-but-not-included labels are supplied as non-load-bearing pointers with their pinned tags where the snapshot tag index has them.\par\medskip
\noindent\textit{Editorial note.} extra\_packages is empty on purpose: the retained preamble already loads xy with options and a second declaration clashes in the assembler wrapper.\par\medskip
\noindent\textit{Editorial note.} Text is retained verbatim, including the author's own typos and placeholders.\par\medskip
\noindent\textit{Editorial note.} Queue index 10 of queue\_examples-defos.json (pb=1885, stmt 2773-2785, proof 2787-2832).\par\medskip
\noindent\textit{Allowed standard prerequisites (editorial).}\begin{itemize}
\item Explicit prerequisite boundary. The cited inputs are prerequisites; the argument in the retained span is not replaced by any cited result.
\end{itemize}\medskip
\ifdefined\hypertarget\hypertarget{sourcepointer-1}{}\fi
\noindent\textit{Source reference.} Original lemma 00S1: sources/stacks-examples-defos/pinned-original-source.tex, source lines 2832--2832. Cited by the retained statement or proof but not part of the collected spans; supplied only as a navigation pointer. The argument does not depend on its content.\par\medskip
\ifdefined\hypertarget\hypertarget{sourcepointer-2}{}\fi
\noindent\textit{Source reference.} Original lemma 00S0: sources/stacks-examples-defos/pinned-original-source.tex, source lines 2832--2832. Cited by the retained statement or proof but not part of the collected spans; supplied only as a navigation pointer. The argument does not depend on its content.\par\medskip
\ifdefined\hypertarget\hypertarget{sourcepointer-3}{}\fi
\noindent\textit{Source reference.} Original lemma 0DZ4: sources/stacks-examples-defos/pinned-original-source.tex, source lines 2383--2393. Cited by the retained statement or proof but not part of the collected spans; supplied only as a navigation pointer. The argument does not depend on its content.\par\medskip
\ifdefined\hypertarget\hypertarget{sourcepointer-4}{}\fi
\noindent\textit{Source reference.} Original lemma 0DYT: sources/stacks-examples-defos/pinned-original-source.tex, source lines 1251--1260. Cited by the retained statement or proof but not part of the collected spans; supplied only as a navigation pointer. The argument does not depend on its content.\par\medskip
\ifdefined\hypertarget\hypertarget{sourcepointer-5}{}\fi
\noindent\textit{Source reference.} Original lemma 0A6A: sources/stacks-examples-defos/pinned-original-source.tex, source lines 2832--2832. Cited by the retained statement or proof but not part of the collected spans; supplied only as a navigation pointer. The argument does not depend on its content.\par\medskip
\ifdefined\hypertarget\hypertarget{sourcepointer-6}{}\fi
\noindent\textit{Source reference.} Original lemma 0D08: sources/stacks-examples-defos/pinned-original-source.tex, source lines 2832--2832. Cited by the retained statement or proof but not part of the collected spans; supplied only as a navigation pointer. The argument does not depend on its content.\par\medskip
\ifdefined\hypertarget\hypertarget{sourcepointer-7}{}\fi
\noindent\textit{Source reference.} Original lemma 0AGU: sources/stacks-examples-defos/pinned-original-source.tex, source lines 2832--2832. Cited by the retained statement or proof but not part of the collected spans; supplied only as a navigation pointer. The argument does not depend on its content.\par\medskip
\ifdefined\hypertarget\hypertarget{sourcepointer-8}{}\fi
\noindent\textit{Source reference.} Original lemma 05EC: sources/stacks-examples-defos/pinned-original-source.tex, source lines 2832--2832. Cited by the retained statement or proof but not part of the collected spans; supplied only as a navigation pointer. The argument does not depend on its content.\par\medskip
\begin{lemma}
\label{lemma-first-order-henselization}
Let $k$ be a field. Let $B$ be a finite type $k$-algebra.
Let $J \subset B$ be an ideal such that
$\Spec(B) \to \Spec(k)$ is smooth on the complement of $V(J)$.
Let $N$ be a finite $B$-module.
Then there is a canonical bijection
$$
\text{Exal}_k(B, N) \to \text{Exal}_k(B^h, N^h)
$$
Here $(B^h, J^h)$ is the henselization of $(B, J)$
and $N^h = N \otimes_B B^h$.
\end{lemma}

\begin{proof}
This proof is the same as the proof of
Lemma \ref{lemma-first-order-completion} but easier. We suggest the
reader to skip the proof.
The map is given by henselization: given $0 \to N \to C \to B \to 0$
in $\text{Exal}_k(B, N)$ we send it to the
henselization $C^h$
of $C$ with respect to the inverse image $J_C \subset C$ of $J$.
Compare with the proof of Lemma \ref{lemma-henselization}.

\medskip\noindent
Since $k \to B$ is of finite presentation the complex
$\NL_{B/k}$ can be represented by a complex
$N^{-1} \to N^0$ where $N^i$ is a finite $B$-module, see
Algebra, Section \ref{algebra-section-netherlander} and
in particular
Algebra, Lemma \ref{algebra-lemma-NL-homotopy}.
As $B$ is Noetherian, this means that $\NL_{B/k}$
is pseudo-coherent. For $g \in J$ the $k$-algebra $B_g$
is smooth and hence $(\NL_{B/k})_g = \NL_{B_g/k}$
is quasi-isomorphic to a finite projective $B$-module sitting in degree $0$.
Thus $\text{Ext}^i_B(\NL_{B/k}, N)_g = 0$ for $i \geq 1$
and any $B$-module $N$. Finally, we have
\begin{align*}
\text{Ext}^1_{B^h}(\NL_{B^h/k}, N^h)
& =
\text{Ext}^1_{B^h}(\NL_{B/k} \otimes_B B^h, N \otimes_B B^h) \\
& =
\text{Ext}^1_B(\NL_{B/k}, N) \otimes_B B^h \\
& =
\text{Ext}^1_B(\NL_{B/k}, N)
\end{align*}
The first equality by
More on Algebra, Lemma \ref{more-algebra-lemma-henselization-NL}
(or rather its analogue for henselizations of pairs).
The second by
More on Algebra, Lemma \ref{more-algebra-lemma-base-change-RHom}.
The third because $\text{Ext}^1_B(\NL_{B/k}, N)$ is $J$-power
torsion, the map $B \to B^h$ is flat and induces an isomorphism
$B/J \to B^h/JB^h$ (More on Algebra, Lemma
\ref{more-algebra-lemma-henselization-flat}), and
More on Algebra, Lemma \ref{more-algebra-lemma-neighbourhood-equivalence}.
This concludes the proof by the description of $\text{Exal}_A(B, N)$
as $\text{Ext}^1_B(\NL_{B/A}, N)$ given just above
Lemma \ref{lemma-first-order-completion}.
\end{proof}

\end{document}
